\problemname{Linguistic Labyrinth}

\illustration{0.28}{illustration.pdf}{
    Spoiler alert: this is one of the~secret test cases.%
    \vspace{-2em}% Ugly negative vspace to make sure the itemize list is not line-wrapped
}

\newcommand{\maxn}{22}
\newcommand{\threeD}{$3$-dimensional }

It is the 25th of December, 2025.
As a Christmas tradition, you gather a group of friends to solve a puzzle.
Among your friends are \emph{wordcels} and \emph{shape rotators},
who are respectively better at thinking with words and with mental images.
This puzzle challenges even the smartest wordcel and the most brilliant shape rotator:

% There is a grid of points in $3$ dimensions, where at each position $(x,y,z)$ with integers $1\leq x,y,z\leq n$ there is a point,
There is a \threeD grid with points at all integer coordinates $(x,y,z)$ with $1\leq x,y,z\leq n$,
and each point has a label associated with it, which is either `\texttt{B}', `\texttt{A}', `\texttt{P}', or `\texttt{C}'.
In this grid, you need to find occurrences of the curly word ``\texttt{BAPC}''.
A curly word ``\texttt{BAPC}'' is a collection of four points in the grid such that:

\begin{itemize}
    \item The labels spell out ``\texttt{BAPC}'' (in this order).
    \item The angle that the triplet ``\texttt{BAP}'' makes is $90$ degrees: the vectors from $\texttt{B}\to \texttt{A}$ and from $\texttt{A}\to \texttt{P}$ form a $90$-degree angle.
    \item The angle that the triplet ``\texttt{APC}'' makes is $90$ degrees: the vectors from $\texttt{A}\to \texttt{P}$ and from $\texttt{P}\to \texttt{C}$ form a $90$-degree angle.
\end{itemize}

Note that the two angles do not need to be axis-aligned.
As an example, see the third sample case, visualized in Figure~\ref{fig:lingu}.

How many occurrences of the curly word ``\texttt{BAPC}'' are in the given grid?

\begin{Input}
    The input consists of:
    \begin{itemize}
        \item One line with an integer $n$ ($1\leq n\leq \maxn$), the size of the grid.
        \item $n$ blocks of $n+1$ lines. Each block of $n+1$ lines consists of:
        \begin{itemize}
            \item One line with a hyphen (\texttt{-}), to make the input more human-readable.
            \item $n$ lines with $n$ characters, each character being either `\texttt{B}', `\texttt{A}', `\texttt{P}', or `\texttt{C}', \\
                representing all labels of one horizontal layer of the \threeD grid.
        \end{itemize}
    \end{itemize}
\end{Input}

\begin{Output}
    Output the number of curly words ``\texttt{BAPC}'' in the \threeD grid.
\end{Output}

\remainingsamples

\vfill

\begin{figure}[h]
    \centering
    \includegraphics[width=.5\textwidth]{sample.pdf}

    \caption{
        Visualization of the third sample input.
        In this grid, there are two curly words ``\texttt{BAPC}'',
        using the highlighted letters.
    }
    \label{fig:lingu}
\end{figure}
